% Work/MAT | 1. rocnik | 22. hodina
% Nove navazuje na diskusni hodinu 21. Definice Pythagorovy vety v obsahove forme.
% Zdroje: ucebnice 03 Planimetrie, 04 Goniometricke funkce,
% Realisticky.cz 030205 a 040201. Texty a TikZ kresby nove sestaveny.
\documentclass[11pt,a4paper]{article}
% Spolecny styl pro ucitelske pripravy MAT (Work/MAT)
% Verze 1.1, 2026-10-23
\NeedsTeXFormat{LaTeX2e}
\ProvidesPackage{priprava}[2026/10/23 v1.1 Sablona ucitelskych priprav MAT]

% Zakladni typografie a cestina
\RequirePackage[T1]{fontenc}
\RequirePackage[utf8]{inputenc}
\RequirePackage[czech]{babel}
\RequirePackage{lmodern}
\RequirePackage{microtype}
\RequirePackage{amsmath,amssymb}
\RequirePackage{xcolor}
\RequirePackage{enumitem}
\RequirePackage[most]{tcolorbox}
\RequirePackage{tabularx,booktabs,array}
\RequirePackage{fancyhdr}
\RequirePackage{needspace}
\RequirePackage{geometry}
\RequirePackage[hidelinks]{hyperref}

\geometry{a4paper,top=17mm,bottom=18mm,left=23mm,right=23mm}
\definecolor{notegray}{gray}{0.39}
\setlength{\parindent}{0pt}
\setlength{\parskip}{2.5pt plus 1pt minus 0.5pt}
\setlist[itemize]{topsep=2pt,itemsep=1pt,parsep=0pt,leftmargin=1.45em}
\setlist[enumerate]{topsep=2pt,itemsep=2pt,parsep=0pt,leftmargin=1.65em}
\setcounter{secnumdepth}{0}
\renewcommand{\arraystretch}{1.12}

% Zahlavi: vlevo rocnik | globalni tema, vpravo cislo: nazev hodiny
% Zapati zustava zcela prazdne.
\pagestyle{fancy}
\fancyhf{}
\renewcommand{\headrulewidth}{0.35pt}
\renewcommand{\footrulewidth}{0pt}
\setlength{\headheight}{14pt}
\newcommand{\nastavPripravu}[4]{%
  \fancyhead[L]{\small #1\enspace\textbar\enspace #2}%
  \fancyhead[R]{\small #3: #4}%
}

% Minimalisticke cernobile boxy. Kazda definice zvlast.
\tcbset{commonlesson/.style={enhanced,breakable=false,colback=white,colframe=black,
  colbacktitle=white,coltitle=black,boxrule=0.45pt,
  arc=0.8mm,outer arc=0.8mm,left=3mm,right=3mm,top=1.6mm,bottom=1.6mm,
  toptitle=1mm,bottomtitle=0.7mm,fonttitle=\bfseries,
  attach boxed title to top left={xshift=2.5mm,yshift=-1.15mm},
  boxed title style={colback=white,colframe=white,boxrule=0pt,left=1mm,right=1mm}}}
\newtcolorbox{zadaniBox}[1][]{commonlesson,title={Zadání},#1}
\newtcolorbox{prikladBox}[1][]{commonlesson,title={Příklad},#1}
\newtcolorbox{definiceBox}[1][]{commonlesson,title={Definice},#1}

% Spolecny box pro zadani a jeho reseni, bez mezery mezi castmi.
% Pouziti:
%   \begin{prikladReseniBox}
%     Zadani ulohy.
%     \tcblower
%     \begin{align*} ... \end{align*}
%   \end{prikladReseniBox}
% Spodni cast ma jemny sedy podklad a delici linku; lze v ni sazet
% vice radku matematiky. Obe casti maji jen jeden spolecny ramecek.
\tcbset{solutionlesson/.style={
  commonlesson,
  enhanced jigsaw,
  breakable,
  colbacklower=black!3,
  segmentation style={solid,draw=black!25,line width=0.35pt},
  fontlower=\small\color{notegray},
  before lower={\noindent\textbf{Řešení:}\enspace},
  leftlower=3mm,rightlower=3mm
}}
\newtcolorbox{prikladReseniBox}[1][]{solutionlesson,title={Příklad},#1}
\newtcolorbox{zadaniReseniBox}[1][]{solutionlesson,title={Zadání},#1}

% Pomucky pro autora pripravy: sedive poznamky; \reseni pro starsi pripravy
\newcommand{\poznamka}[1]{\par\begingroup\small\color{notegray}\itshape
  \textnormal{\textbf{Poznámka pro učitele.}}\enspace #1\par\endgroup}
\newcommand{\reseni}[1]{\par\begingroup\small\color{notegray}
  \textbf{Očekávané řešení:}\enspace #1\par\endgroup}
\newcommand{\blok}[1]{\needspace{5\baselineskip}\subsection*{#1}}
\newcommand{\vyrok}[1]{\uv{#1}}

\usepackage{tikz}
\usetikzlibrary{calc}
\nastavPripravu{1. ročník}{Goniometrie}{22}{Trojúhelník, Pythagoras, poměry}
\begin{document}
\blok{Od debaty k pravoúhlému trojúhelníku (0.--7. minuta)}
Připomeň včerejší dvě rampy: stejný \emph{sklon} může mít různě velký obrázek.
V geometrii nás bude zajímat úhel a poměr správně vybraných stran.
\begin{definiceBox}
\textbf{Pravoúhlý trojúhelník} je trojúhelník s jedním vnitřním úhlem
$90^\circ$. Zbývající dva úhly jsou ostré a jejich součet je $90^\circ$.
\end{definiceBox}
\begin{definiceBox}
\textbf{Přepona} je strana proti pravému úhlu. \textbf{Odvěsny} jsou
zbývající dvě strany, které pravý úhel svírají.
\end{definiceBox}
\begin{center}
\begin{tikzpicture}[scale=.78]
\coordinate (A) at (0,0); \coordinate (C) at (5.2,0); \coordinate (B) at (5.2,3.15);
\draw[very thick] (A)--node[below]{$b$}(C)
 --node[right]{$a$}(B)
 --node[above,sloped]{$c$}(A);
\draw (4.88,0)--(4.88,.32)--(5.2,.32);
\draw (.84,0) arc[start angle=0,end angle=31.2,radius=.84];
\node at (1.15,.28) {$\alpha$};
\node[below left] at (A){$A$};\node[below right] at (C){$C$};\node[above right] at (B){$B$};
\node at (5.65,.2){$90^\circ$};
\end{tikzpicture}
\end{center}
Pro úhel $\alpha$ je odvěsna $a$ \emph{protilehlá}, $b$ \emph{přilehlá} a $c$ přepona.
\poznamka{V obrázku je $\gamma=90^\circ$, tedy přeponou je $c$; ostatní
popisky stran odpovídají obvyklému značení stran proti vrcholům.
V jiném trojúhelníku může být přeponou $a$ nebo $b$.
Nedovol žákům rozhodovat podle toho, která čára je šikmá.}

\blok{Pythagorova věta vyjadřuje rovnost obsahů (7.--15. minuta)}
\begin{definiceBox}
\textbf{Pythagorova věta:} \emph{Obsah čtverce nad přeponou
pravoúhlého trojúhelníku je roven součtu obsahů čtverců nad oběma
odvěsnami.} Jsou-li odvěsny $a,b$ a přepona $c$, platí
\[
\boxed{a^2+b^2=c^2}.
\]
Zápis $c^2=a^2+b^2$ používáme jen tehdy, je-li $c$ skutečně přepona.
\end{definiceBox}
\begin{center}
\begin{tikzpicture}[scale=.53]
\draw[thick] (0,0) rectangle (5,5);
\fill[black!12] (3,0)--(5,3)--(2,5)--(0,2)--cycle;
\draw[very thick] (3,0)--(5,3)--(2,5)--(0,2)--cycle;
\node at (2.5,2.5) {$c^2$};
\node[below] at (1.5,0) {$a$}; \node[below] at (4,0) {$b$};
\node[left] at (0,2.5) {$a+b$};
\end{tikzpicture}
\end{center}
Srovnej obsah velkého čtverce: lze jej rozdělit na čtverce
$a^2$, $b^2$ a dva obdélníky $ab$, nebo na čtyři trojúhelníky
o obsahu $\frac{ab}{2}$ a prostřední čtverec $c^2$.
Proto $a^2+2ab+b^2=2ab+c^2$, tedy $a^2+b^2=c^2$.
\poznamka{Důkaz je názorný. Vnitřní čtyřúhelník je čtverec,
protože dva sousední ostré úhly dohromady tvoří $90^\circ$.}
\poznamka{Časový rámec: strany 7, Pythagoras 8,
čtyři poměry a podobnost 13, P1--P2 a S1 14, závěr 3 minut.
Domácí po této hodině: u trojúhelníku $8$--$15$--$17$ napiš
čtyři poměry pro úhel proti osmičce. Výsledky
$8/17$, $15/17$, $8/15$, $15/8$.
Speciální hodnoty pro konkrétní úhly až příště.}

\blok{Poměr závisí na úhlu, ne na velikosti obrázku (15.--28. minuta)}
Dva pravoúhlé trojúhelníky se stejným ostrým úhlem jsou podobné
(podle věty \emph{uu}); při zvětšení se všechny délky násobí stejným
číslem. \textbf{Poměr odpovídajících délek se nezmění.}
Vzhledem ke zvolenému ostrému úhlu $\alpha$ určujeme \emph{protilehlou}
a \emph{přilehlou odvěsnu}; přepona je vždy tatáž strana.
\begin{definiceBox}
\textbf{Sinus} ostrého úhlu je poměr protilehlé odvěsny a přepony:
$\displaystyle \sin\alpha=\frac{\text{protilehlá odvěsna}}{\text{přepona}}$.
\end{definiceBox}
\begin{definiceBox}
\textbf{Kosinus} ostrého úhlu je poměr přilehlé odvěsny a přepony:
$\displaystyle \cos\alpha=\frac{\text{přilehlá odvěsna}}{\text{přepona}}$.
\end{definiceBox}
\begin{definiceBox}
\textbf{Tangens} ostrého úhlu je poměr protilehlé a přilehlé odvěsny:
$\displaystyle \tan\alpha=\frac{\text{protilehlá odvěsna}}{\text{přilehlá odvěsna}}$.
Značí se také $\operatorname{tg}\alpha$.
\end{definiceBox}
\begin{definiceBox}
\textbf{Kotangens} ostrého úhlu je poměr přilehlé a protilehlé odvěsny:
$\displaystyle \cot\alpha=\frac{\text{přilehlá odvěsna}}{\text{protilehlá odvěsna}}$.
Značí se také $\operatorname{cotg}\alpha$.
\end{definiceBox}
\textbf{V tomto modelu} definujeme poměry pouze pro ostré úhly,
$0^\circ<\alpha<90^\circ$. Hodnotám na krajích se budeme věnovat zítra.

\begin{prikladReseniBox}[title={P1: trojúhelník $3$--$4$--$5$},breakable=false]
Pravoúhlý trojúhelník má vzhledem k úhlu $\alpha$ protilehlou odvěsnu $3\,\mathrm{cm}$,
přilehlou $4\,\mathrm{cm}$ a přeponu $5\,\mathrm{cm}$.
Vyjádři čtyři goniometrické poměry. Co se změní pro druhý ostrý úhel $\beta$?
\tcblower
$\sin\alpha=\frac35$, $\cos\alpha=\frac45$, $\tan\alpha=\frac34$,
$\cot\alpha=\frac43$. Pro $\beta$ si odvěsny prohodí role:
$\sin\beta=\frac45$, $\cos\beta=\frac35$, $\tan\beta=\frac43$,
$\cot\beta=\frac34$. V každém případě jsou poměry bez jednotek.
\end{prikladReseniBox}

\blok{Zkus určit přeponu a poměry (28.--42. minuta)}
\begin{prikladReseniBox}[title={P2: nejprve Pythagorova věta},breakable=false]
Odvěsny jsou $5\,\mathrm{cm}$ a $12\,\mathrm{cm}$. Urči přeponu.
Pak zapiš všechny čtyři poměry pro úhel proti kratší odvěsně.
\tcblower
$c=\sqrt{5^2+12^2}=\sqrt{169}=13\,\mathrm{cm}$.
$\sin\alpha=\frac5{13}$, $\cos\alpha=\frac{12}{13}$,
$\tan\alpha=\frac5{12}$, $\cot\alpha=\frac{12}5$.
Kontrola: přepona je nejdelší a $5^2+12^2=13^2$.
\end{prikladReseniBox}
\begin{zadaniReseniBox}[title={S1: samostatně (jádro a)--c))},breakable=false]
\begin{enumerate}[label=\alph*)]
\item Odvěsny měří $6\,\mathrm{cm}$ a $8\,\mathrm{cm}$. Urči přeponu.
\item Pro odvěsny z a) zapiš $\sin,\cos,\tan,\cot$ úhlu proti odvěsně $6\,\mathrm{cm}$.
\item Je-li pravý úhel $\beta=90^\circ$, která ze stran $a,b,c$ je přeponou?
\item (Navíc) Vysvětli, proč se hodnoty poměrů nezmění, když všechny strany zdvojnásobíme.
\end{enumerate}
\tcblower
a) $\sqrt{36+64}=10\,\mathrm{cm}$.
b) $\frac35,\frac45,\frac34,\frac43$.
c) Strana $b$, protože leží proti pravému úhlu.
d) Každý poměr má v čitateli i jmenovateli stejný násobek $2$, který se zkrátí.
\end{zadaniReseniBox}
\end{document}
